%%%%%%%%%%%%%%%%%%%%%part4.tex%%%%%%%%%%%%%%%%%%%%%%%%%%%%%%%%%%
% 
% Part IV: Model Robustness & Transparency
%
%%%%%%%%%%%%%%%%%%%%%%%% Springer Nature %%%%%%%%%%%%%%%%%%%%%%

\begin{partbacktext}
\part{Model Robustness \& Transparency}
\noindent To be truly trustworthy, an AI system must not only be private and secure but also robust to manipulation and transparent in its decision-making. This fourth part of the book shifts the focus from defending the data to hardening the model itself.

We begin in \textbf{Chapter 11} with \textbf{Adversarial Hardening}, exploring how to defend models against evasion attacks. In \textbf{Chapter 12}, we build an \textbf{Explainability Stack} (XAI) using tools like SHAP to ensure model outputs are interpretable. \textbf{Chapter 13} addresses the critical dimension of \textbf{Algorithmic Fairness}, exploring how to measure and mitigate demographic bias. \textbf{Chapter 14} addresses the emerging threat of \textbf{Post-Quantum Privacy}, securing our cryptographic foundations for the future. \textbf{Chapter 15} explores \textbf{Efficient AI}, demonstrating how to make privacy-preserving techniques practical through quantization and pruning. Finally, \textbf{Chapter 16} introduces the critical domain of \textbf{Intellectual Property and Copyright}, exploring cryptographic watermarking and provenance in GenAI.

\begin{casestudy} 
The 2025 McHire scandal began as a minor calibration error in an automated recruitment agent. Over six months, the model"s feedback loop amplified a slight preference for candidates from specific universities into a systemic exclusion of 92 percent of minority applicants. This incident proved that a technically "secure" model can still be socially catastrophic, leading to the rigorous Fairness, Robustness, and Transparency frameworks of Part IV. 
\end{casestudy} 
\end{partbacktext}
